\chapter{Adjektive}
\label{cha:Adjektive}

Adjektive beschreiben oder bestimmen Nomen näher.

\section{Deklination der Adjektive}
\label{sec:AdjDeklination}

Adjektive werden dekliniert und richten sich nach Genus, Numerus und Kasus des bezogenen Nomens.

\subsection{Starke Deklination (ohne Artikel)}
\begin{center}
\begin{tabular}{|l|c|c|c|c|}
\hline
\textbf{Kasus} & \textbf{Maskulin} & \textbf{Feminin} & \textbf{Neutral} & \textbf{Plural} \\ \hline
Nominativ & gut-er & gut-e & gut-es & gut-e \\ \hline
Akkusativ & gut-en & gut-e & gut-es & gut-e \\ \hline
Dativ & gut-em & gut-er & gut-em & gut-en \\ \hline
Genitiv & gut-en & gut-er & gut-en & gut-er \\ \hline
\end{tabular}
\end{center}

\subsection{Gemischte Deklination (unbestimmter Artikel)}
\begin{center}
\begin{tabular}{|l|c|c|c|c|}
\hline
\textbf{Kasus} & \textbf{Maskulin} & \textbf{Feminin} & \textbf{Neutral} & \textbf{Plural} \\ \hline
Nominativ & gut-er & gut-e & gut-es & gut-en \\ \hline
Akkusativ & gut-en & gut-e & gut-es & gut-en \\ \hline
Dativ & gut-en & gut-er & gut-en & gut-en \\ \hline
Genitiv & gut-en & gut-er & gut-en & gut-en \\ \hline
\end{tabular}
\end{center}

\subsection{Schwache Deklination (bestimmter Artikel)}
\begin{center}
\begin{tabular}{|l|c|c|c|c|}
\hline
\textbf{Kasus} & \textbf{Maskulin} & \textbf{Feminin} & \textbf{Neutral} & \textbf{Plural} \\ \hline
Nominativ & gut-e & gut-e & gut-e & gut-en \\ \hline
Akkusativ & gut-en & gut-e & gut-e & gut-en \\ \hline
Dativ & gut-en & gut-en & gut-en & gut-en \\ \hline
Genitiv & gut-en & gut-en & gut-en & gut-en \\ \hline
\end{tabular}
\end{center}

\section{Steigerung (Komparativ, Superlativ)}
\label{sec:AdjSteigerung}

\subsection{Regelmäßige Steigerung}
\begin{itemize}
    \item \textbf{Positiv}: klein
    \item \textbf{Komparativ}: kleiner
    \item \textbf{Superlativ}: am kleinsten
\end{itemize}

\subsection{Unregelmäßige Steigerung}
\begin{center}
\begin{tabular}{|l|c|c|}
\hline
\textbf{Positiv} & \textbf{Komparativ} & \textbf{Superlativ} \\ \hline
gut & besser & am besten \\ \hline
viel & mehr & am meisten \\ \hline
gern & lieber & am liebsten \\ \hline
groß & größer & am größten \\ \hline
hoch & höher & am höchsten \\ \hline
nah & näher & am nächsten \\ \hline
\end{tabular}
\end{center}

\section{Adjektiv vs. Adverb}
\label{sec:AdjvsAdv}

\begin{tcolorbox}[title=Unterscheidung]
\begin{itemize}
    \item \textbf{Adjektiv}: beschreibt das Nomen
    \item \textbf{Adverb}: beschreibt das Verb
\end{itemize}
\end{tcolorbox}

\begin{itemize}
    \item Er ist ein \textbf{guter} Schüler. (Adjektiv)
    \item Er lernt \textbf{gut}. (Adverb)
\end{itemize}

\section{Übungen zu Adjektiven}
\label{sec:AdjUebungen}

\begin{enumerate}
    \item Das ist ein \underline{\hspace{1cm}} Buch. (interessant)
    \item Er hat ein \underline{\hspace{1cm}} Auto. (neu)
    \item Die \underline{\hspace{1cm}} Frau ist meine Mutter. (alt)
    \item Ich habe \underline{\hspace{1cm}} Probleme. (viel)
    \item Das ist der \underline{\hspace{1cm}} Weg. (kurz)
\end{enumerate}

\chapterend